\section{User Interface and Popup Menu Guide}
\label{sec:ui-popup-menu}

\FractalShark{} can be driven \emph{entirely} from the keyboard and mouse, with the
right-click popup menu acting as the primary discovery and configuration
surface.  Day-to-day navigation (zooming, recentering, going back, iteration
adjustment, palette cycling, and several recomputation paths) is also exposed
through single-key shortcuts, so common workflows do not require opening the
menu.

The popup menu is context-sensitive and rebuilt dynamically each time it is
shown, so that checked, enabled, and disabled states accurately reflect the
current application settings and rendering state.  The menu can be invoked by
mouse (\textbf{right-click}) or from the keyboard via the standard Windows
context-menu invocation (\textbf{Shift+F10} / \textbf{Menu} key).

\subsection{Basic Interaction}

On startup, \FractalShark{} displays a brief splash screen (randomly chosen
from a set of built-in images) while initializing the renderer.

The most common interactions are:
\begin{itemize}
  \item \textbf{Right-click}: Open the popup menu at the mouse cursor.
  \item \textbf{Left-click and drag}: Zoom into a rectangular region.
        By default the zoom rectangle preserves the window aspect ratio;
        holding \code{Shift} while dragging allows a free-form
        (non-square) selection.
  \item \textbf{Mouse wheel}: Scroll forward to zoom in at the cursor;
        scroll backward to zoom out from the center.
  \item \textbf{Keyboard shortcuts}:
    \begin{itemize}
      \item \code{z}: Zoom in at the mouse cursor.
      \item \code{Shift+Z}: Zoom out slightly.
      \item \code{b}: Return to the previous view.
      \item \code{- / =}: Decrease or increase the iteration limit.
    \end{itemize}
\end{itemize}

Resizing the window automatically resets the renderer dimensions and triggers
a recalculation of the fractal to match the new window size.

\subsection{Menu Structure and Behavior}

The popup menu is organized into logical groups and submenus. Items fall into
three distinct behavioral categories:

\begin{description}
  \item[Radio groups] Mutually exclusive choices where exactly one option is
  active at a time (for example, selecting a render algorithm or GPU antialiasing level).
  \item[Toggles] Binary on/off options that switch between two states each time
  they are selected (for example, toggling repainting or full-screen mode).
  \item[Commands] Immediate actions that execute when selected and do not
  maintain a checked state.
\end{description}

Radio groups reflect the configuration in use, with the currently active option
checked. Selecting a different option automatically deselects the previous
choice and triggers any necessary reconfiguration.

\subsection{Top-Level Menu}

\begin{itemize}
  \item \emph{Show Help} (command): Displays an on-screen help dialog summarizing keyboard shortcuts.

  \item \emph{Navigate} (submenu):
    \begin{itemize}
      \item \emph{Back} (command): Navigate back to the previous view in history.
      \item \emph{Center View Here} (command): Centers the current view on the
      point under the cursor.
      \item \emph{Zoom In Here} (command): Zooms in toward the point under the cursor.
      \item \emph{Zoom Out} (command): Zooms out from the current view.
      \item \emph{Autozoom Max} (command): Autozoom using the ``Max'' heuristic,
      which zooms by a factor of~32 each step.  See \cref{par:autozoom} for details.
      \item \emph{Autozoom Default} (command): See \cref{par:autozoom} for details.
      \item \emph{Autozoom Filament Tip} (\code{S}, i.e.\ \code{Shift+s}) (command): Autozoom using the
      ``FilamentTip'' heuristic.
      \item \emph{Feature Finder} (submenu): Locates nearby mini-Mandelbrot
      nuclei (periodic points) using Newton's method and optionally zooms to
      them.  Three evaluation back-ends are available, each in a single-point
      and a whole-screen scan variant.  See
      \cref{sec:feature-finder} for the underlying algorithm.
        \begin{itemize}
          \item \emph{Direct} (\code{n}) (command): Runs the feature finder at
          the mouse cursor using direct iteration.  Best for shallow zooms where
          no perturbation data is needed.
          \item \emph{DirectScan} (\code{N}) (command): Runs direct-iteration
          feature finding on a $12\times12$ grid across the entire viewport,
          locating multiple features at once.
          \item \emph{PT} (\code{m}) (command): Runs the feature finder at the
          mouse cursor using perturbation theory (\cref{sec:perturbation-concept}).
          Requires a reference orbit and works at deeper zooms than Direct.
          \item \emph{PTScan} (\code{M}) (command): Runs perturbation-theory
          feature finding on a $12\times12$ grid across the viewport.
          \item \emph{LA} (\code{,}) (command): Runs the feature finder at the
          mouse cursor using linear approximation (\cref{sec:la-v2-perturb}).
          Works at extreme zoom depths where PT alone would be too slow.
          \item \emph{LAScan} (\code{<}) (command): Runs linear-approximation
          feature finding on a $12\times12$ grid across the viewport.
          \item \emph{Zoom to Found Feature} (\code{.}) (command): Zooms into
          the most recently found feature.  If the feature has only been located
          at low precision, a high-precision GMP refinement pass is performed
          automatically before zooming.
          \item \emph{Clear All Found Features} (\code{>}) (command): Removes
          all found-feature overlays (lines and circles drawn on the fractal).
          \item \emph{Resume NR Refinement} (command): Resumes
          Newton--Raphson precision refinement from the most recent checkpoint,
          useful when a prior refinement was interrupted or when additional
          precision is desired.
          \item \emph{NR Inner Loop: GPU / CPU} (radio): Selects
          whether the Newton--Raphson inner loop (iterating the orbit and
          derivatives for $p$~steps) executes on the GPU via
          \code{HpSharkFloat} or falls back to CPU-based MPIR arithmetic.
        \end{itemize}
      When a feature is found, a line is drawn from the search origin to the
      detected nucleus, making its location visible.  Scan modes are useful for
      surveying an entire view at a glance; single-point modes are faster and
      target the area near the cursor.
    \end{itemize}

  \item \emph{Built-In Views} (submenu): Loads a predefined view/location used
  for demonstrations, debugging, or performance testing.  Up to 32 views are
  defined in the source (\code{FractalViewPresets.cpp}).
    \begin{itemize}
      \item \emph{Help} (command): Pops up a message box that simply says this
      submenu contains built-in test views.
      \item \emph{Standard View} (command): Loads the default starting view: the
      classic Mandelbrot.
      \item \emph{\#1 -- \#32} (commands): Loads a specific predefined
      test view. Many entries encode expected precision limits, kernel behavior,
      or known failure cases in their label text.  Several views require intense
      computation and take a long time to render.  View \#27, for example, is a
      period ~28 billion point that takes hours to render and requires reference
      compression even with 128GB of RAM.  Exercise caution when selecting deep-zoom
      views, as they may cause long delays or high memory usage.

    \end{itemize}
    When reference orbits or Imagina-format orbit files are loaded, additional
    dynamic view entries may appear in the menu below the built-in presets.
    These entries correspond to the loaded data and are removed when the orbit
    state is cleared.

  \item \emph{Recalculate, Reuse Reference} (command): Forces an immediate
  re-render using the current settings, reusing the existing reference orbit if
  possible.  See \cref{subsec:recalculation-commands} for details.

  \item \emph{Toggle Repainting}: Enables or disables progressive
  repainting while a render is in progress. When disabled, no recalculation
  takes place on window resize or related, which can be helpful when a new
  window size is desired without triggering a re-render.
  \item \emph{Toggle Window Size}: Toggles between windowed and
  full-screen behavior.  The window stays at the same aspect ratio by default.
  \item \emph{Toggle Window Size (Square)}: Toggles an alternate sizing
  mode constrained to a square aspect ratio.
  \item \emph{Minimize Window} (command): Minimizes the application window.

  \item \emph{Antialiasing} (radio group submenu): Selects the antialiasing
  level used for rendering.
    \begin{itemize}
      \item \emph{1x (fast)}
      \item \emph{4x}
      \item \emph{9x}
      \item \emph{16x (better quality)}
    \end{itemize}

  Note that antialiasing as used here is not traditional GPU antialiasing.
  Instead, 4, 9, or 16 samples are explicitly computed per pixel and averaged
  together either on the CPU or GPU.  This approach is known as
  \emph{supersampling} and produces high-quality results at the cost of
  increased computation time and memory usage.  It's a naive approach---in
  principle, more sophisticated techniques could be implemented in the future,
  such as edge-aware sampling or adaptive sampling.

  \item \emph{Choose Render Algorithm} (radio group submenu): Selects the core
  Mandelbrot rendering strategy. This submenu is a radio group: selecting one algorithm
  automatically deselects the others. Available choices may include CPU-based
  reference and linear-approximation renderers, multiple CUDA-based GPU
  renderers, and high-dynamic-range and experimental kernels. On systems without
  a working CUDA setup, CPU renderers can be selected as a fallback. In the
  event of a CUDA initialization failure, switching to a CPU algorithm followed
  by a recalculation may restore a usable image.

  Algorithm names follow a systematic convention:
  \begin{description}
    \item[\textrm{N}x\textrm{BB}] Precision type: \emph{N} is the number of
    components and \emph{BB} is the component bit-width.
    \code{1x32}~=~single float, \code{2x32}~=~double-float
    (\cref{subsec:ep-dblflt}), \code{1x64}~=~double,
    \code{2x64}~=~double-double, \code{4x32}/\code{4x64}~=~quad
    types (\cref{subsec:ep-higher-prec}).
    \item[HDR] High Dynamic Range: the type is wrapped in
    \code{HDRFloat} (\cref{sec:hdr-float}) for extended exponent range,
    required at deep zooms.
    \item[RC] Reference Compression: the reference orbit is stored in
    compressed form (\cref{sec:ref-compression-zhuoran}).
    \item[Perturbation / Perturb] Per-pixel perturbation from a
    reference orbit (\cref{sec:perturbation-concept}).
    \item[BLA] Bilinear Approximation v1 (\cref{sec:bilinear-approx}).
    \item[LA~v2] Linear Approximation v2 (\cref{sec:la-v2-perturb}),
    the primary deep-zoom rendering path.
    \item[Scaled] An alternative perturbation formulation using a
    rescaled delta; partially broken.
    \item[LA only / Perturb only] Testing modes that run only one
    stage of the pipeline in isolation.
  \end{description}
  For example, \emph{HDRx2x32 GPU -- RC LA~v2} means: GPU kernel using
  a 2$\times$32-bit (double-float) type wrapped in HDR, with reference
  compression and the LA~v2 full pipeline.  Individual algorithm entries
  below are not further annotated because their names are
  self-documenting given this convention.

    \begin{itemize}
      \item \emph{Help} (command): Displays help for algorithm selection and
      related concepts (e.g., what LA~v2 (\cref{sec:la-v2-perturb}), BLA
      (\cref{sec:bilinear-approx}), perturbation (\cref{sec:perturbation-concept}),
      and reference compression (\cref{subsec:ref-orbit-compression-reuse}) mean,
      and when to use each).

      \item \emph{Auto (Default)} (radio): Automatically selects the
      default/recommended renderer.  Recommended.

      \emph{Note}: this option exhibits potentially-confusing behavior.  When
      selected, it does not lock in a specific algorithm; rather, it allows the
      system to choose the most appropriate algorithm based on current view
      parameters. The radio button for "Auto" will not be checked; instead, the
      actual algorithm in use will be checked.  Thus, selecting \emph{Auto} may
      appear to have no effect, but navigating the submenus reveals that the
      algorithm currently in use is checked.  This design allows users to see
      what algorithm is active.

      At present, four algorithms are eligible for automatic selection:
      \begin{itemize}
        \item \emph{1x32 GPU}: The default low-zoom GPU linear-approximation
              renderer.  No perturbation or linear approximation is used.
        \item \emph{1x32 GPU -- Perturbation Only}: The default
              perturbation-only GPU renderer (\cref{sec:perturb-only}) for
              moderate zooms uses perturbation but no linear approximation.
        \item \emph{1x32 GPU -- LA~v2}: A bit deeper, no high-dynamic range
              floats but does use linear-approximation renderer and
              perturbation.
        \item \emph{HDRx32 GPU -- LA~v2}: Arbitrary depth (\cref{sec:la-v2-perturb,sec:hdr-float}), combines all known
              algorithmic optimizations.  This one can result in rendering
              artifacts at some deep locations due to numeric precision
              limitations.  ``Auto'' mode cannot detect this, so if artifacts appear, try
              HDRx2x32 or HDRx64 variants instead.  View \#26 is an example of a
              point that fails with HDRx32 but works correctly with HDRx64.

              See \cref{fig:auto-renderer-artifact} for an example of
              artifacts that can arise when using HDRx32 at extreme depths.
      \end{itemize}

      \item \emph{LA Parameters} (submenu): Controls linear-approximation (LA) configuration and tuning.
        \begin{itemize}
          \item \emph{Multithreaded (Default)} (radio): Uses multithreaded LA processing.
          \item \emph{Single Threaded} (radio): Uses single-threaded LA processing.
          \item \emph{Max Accuracy (Default)} (command): Resets all LA threshold
          and period-detection exponents to their conservative defaults.  Produces the
          most accurate linear approximation at the cost of longer computation.
          \item \emph{Max Performance (Accuracy Loss)} (command): Increases the LA
          threshold scale exponents by~12 relative to defaults, allowing the
          approximation to accept looser tolerances.  This speeds up LA computation
          but may introduce visible artifacts at some locations.
          \item \emph{Min Memory} (command): Increases the period-detection threshold
          exponents by~3, causing the period detector to accept matches more eagerly.
          This setting reduces the number of stored LA stages and lowers memory consumption.
        \end{itemize}

      \item \emph{CPU-Only} (submenu): CPU renderers and CPU-driven perturbation variants.
        \begin{itemize}
          \item \emph{Very High Precision CPU} (radio)
          \item \emph{1x64 CPU} (radio)
          \item \emph{HDRx32 CPU} (radio)
          \item \emph{HDRx64 CPU} (radio)

          \item \emph{1x64 CPU -- Perturbation BLA} (radio)
          \item \emph{HDRx32 CPU -- Perturbation BLA} (radio)
          \item \emph{HDRx64 CPU -- Perturbation BLA} (radio)

          \item \emph{HDRx32 CPU -- Perturbation LA~v2} (radio)
          \item \emph{HDRx64 CPU -- Perturbation LA~v2} (radio)
          \item \emph{HDRx32 RC CPU -- Perturbation LA~v2} (radio)
          \item \emph{HDRx64 RC CPU -- Perturbation LA~v2} (radio)
        \end{itemize}

      \item \emph{Low-Zoom Depth} (submenu): GPU renderers intended for relatively modest zoom depths and lower precision requirements.
        \begin{itemize}
          \item \emph{Iteration Precision} (radio group submenu): Controls a
          rendering speed/quality trade-off for low-zoom GPU paths.  The
          internal multiplier (1, 4, 8, or~16) scales the iteration step size:
          higher multipliers compute faster but with lower fidelity.
            \begin{itemize}
              \item \emph{4x (fast)} (radio): Multiplier~16 (fastest, lowest quality).
              \item \emph{3x} (radio): Multiplier~8.
              \item \emph{2x} (radio): Multiplier~4.
              \item \emph{1x (better quality)} (radio): Multiplier~1 (slowest, highest quality).
            \end{itemize}

          \item \emph{1x32 GPU} (radio)
          \item \emph{2x32 GPU} (radio)
          \item \emph{4x32 GPU} (radio)
          \item \emph{1x64 GPU} (radio)
          \item \emph{2x64 GPU} (radio)
          \item \emph{4x64 GPU} (radio)
          \item \emph{HDRx32 GPU} (radio)
        \end{itemize}

      \item \emph{Scaled} (submenu): Perturbation renderers using a scaled formulation.
        \begin{itemize}
          \item \emph{1x32 GPU -- Perturbation Scaled} (radio)
          \item \emph{2x32 GPU -- Perturbation Scaled (broken)} (radio)
          \item \emph{HDRx32 GPU -- Perturbation Scaled} (radio)
        \end{itemize}

      \item \emph{Bilinear Approximation V1} (submenu): Perturbation renderers using the original bilinear approximation approach.
        \begin{itemize}
          \item \emph{1x64 GPU -- Perturbation BLA} (radio)
          \item \emph{HDRx32 GPU -- Perturbation BLA} (radio)
          \item \emph{HDRx64 GPU -- Perturbation BLA} (radio)
        \end{itemize}

      \item \emph{LA Only (for testing)} (submenu): Test configurations that run only the linear-approximation portion of the pipeline.
        \begin{itemize}
          \item \emph{1x32 GPU -- LA~v2 -- LA only} (radio)
          \item \emph{2x32 GPU -- LA~v2 -- LA only} (radio)
          \item \emph{1x64 GPU -- LA~v2 -- LA only} (radio)
          \item \emph{HDRx32 GPU -- LA~v2 -- LA only} (radio)
          \item \emph{HDRx2x32 GPU -- LA~v2 -- LA only} (radio)
          \item \emph{HDRx64 GPU -- LA~v2 -- LA only} (radio)
        \end{itemize}

      \item \emph{Perturbation Only} (submenu): Test configurations that run only the perturbation portion of the pipeline.
        \begin{itemize}
          \item \emph{1x32 GPU -- Perturb only} (radio)
          \item \emph{2x32 GPU -- Perturb only} (radio)
          \item \emph{1x64 GPU -- Perturb only} (radio)
          \item \emph{HDRx32 GPU -- Perturb only} (radio)
          \item \emph{HDRx2x32 GPU -- Perturb only} (radio)
          \item \emph{HDRx64 GPU -- Perturb only} (radio)
        \end{itemize}

      \item \emph{Reference Compression} (submenu): Algorithms using reference-orbit compression as part of the perturbation pipeline.
        \begin{itemize}
          \item \emph{1x32 GPU -- RC LA~v2} (radio)
          \item \emph{2x32 GPU -- RC LA~v2} (radio)
          \item \emph{1x64 GPU -- RC LA~v2} (radio)
          \item \emph{HDRx32 GPU -- RC LA~v2} (radio)
          \item \emph{HDRx2x32 GPU -- RC LA~v2} (radio)
          \item \emph{HDRx64 GPU -- RC LA~v2} (radio)

          \item \emph{Perturbation Only} (submenu): Reference-compressed perturbation-only variants.
            \begin{itemize}
              \item \emph{1x32 GPU -- RC Perturb Only} (radio)
              \item \emph{2x32 GPU -- RC Perturb Only} (radio)
              \item \emph{1x64 GPU -- RC Perturb Only} (radio)
              \item \emph{HDRx32 GPU -- RC Perturb Only} (radio)
              \item \emph{HDRx2x32 GPU -- RC Perturb Only} (radio)
              \item \emph{HDRx64 GPU -- RC Perturb Only} (radio)
            \end{itemize}

          \item \emph{LA Only} (submenu): Reference-compressed LA-only variants.
            \begin{itemize}
              \item \emph{1x32 GPU -- RC LA~v2 -- LA only} (radio)
              \item \emph{2x32 GPU -- RC LA~v2 -- LA only} (radio)
              \item \emph{1x64 GPU -- RC LA~v2 -- LA only} (radio)
              \item \emph{HDRx32 GPU -- RC LA~v2 -- LA only} (radio)
              \item \emph{HDRx2x32 GPU -- RC LA~v2 -- LA only} (radio)
              \item \emph{HDRx64 GPU -- RC LA~v2 -- LA only} (radio)
            \end{itemize}
        \end{itemize}

      \item \emph{LA~v2 (Full Pipeline)} (radio entries): Main LA~v2 renderers combining LA, perturbation, and reference handling.
        \begin{itemize}
          \item \emph{1x32 GPU -- LA~v2} (radio)
          \item \emph{2x32 GPU -- LA~v2} (radio)
          \item \emph{1x64 GPU -- LA~v2} (radio)
          \item \emph{HDRx32 GPU -- LA~v2} (radio)
          \item \emph{HDRx2x32 GPU -- LA~v2} (radio)
          \item \emph{HDRx64 GPU -- LA~v2} (radio)
        \end{itemize}

      \item \emph{Tests} (submenu): Development and diagnostic test actions.
        \begin{itemize}
          \item \emph{Run Basic Test (saves files in local dir)} (command)
          \item \emph{Run View \#27 test} (command)
        \end{itemize}

    \end{itemize}

  \item \emph{Iterations} (submenu): Adjusts the iteration limit and iteration counter width.
    \begin{itemize}
      \item \emph{Default Iterations} (command): Resets iteration count to defaults.
      \item \emph{+1.5x} (command): Multiplies iteration limit by 1.5.
      \item \emph{+6x} (command): Multiplies iteration limit by 6.
      \item \emph{+24x} (command): Multiplies iteration limit by 24.
      \item \emph{Decrease Iterations} (command): Decreases the current iteration limit.
      \item \emph{32-Bit Iterations} (radio): Uses 32-bit iteration counters (lower memory, smaller range).
      \item \emph{64-Bit Iterations} (radio): Uses 64-bit iteration counters (larger range).
    \end{itemize}

  \item \emph{Perturbation} (radio group submenu): Controls perturbation
    rendering and related reference-orbit management.  \FractalShark{} supports
    perturbation-based rendering driven by a high-precision reference orbit.
    Menu options in this area control how the reference orbit is computed and
    reused.  GPU-accelerated reference orbits are intended for extremely deep
    zooms and may be slower than CPU computation at modest precision levels.
    They are marked as experimental (\cref{subsec:experimental-debug}) and
    should be used accordingly.
    \begin{itemize}
      \item \emph{Clear Perturbation References -- All} (command): Clears all cached perturbation references.
      \item \emph{Clear Perturbation References -- Med} (command): Clears medium-class cached references.
      \item \emph{Clear Perturbation References -- High} (command): Clears high-class cached references.
      \item \emph{Show Perturbation Results} (command): Intended to display
      diagnostic overlay information for perturbation state.  Currently
      a stub (prints a TODO to the debug console).

      \item \emph{Auto (default)} (radio): Automatically selects the perturbation
      algorithm based on current conditions.  Recommended for most use.
      \item \emph{Single Thread (ST)} (radio): Uses a single-threaded
      reference-orbit computation with no periodicity checking.
      \item \emph{Multi Thread (MT)} (radio): Uses a multithreaded
      reference-orbit computation with no periodicity checking.
      \item \emph{ST + Periodicity} (radio): Single-threaded reference-orbit
      computation with periodicity checking to detect convergence early.
      \item \emph{MT2 + Periodicity} (radio): Multithreaded reference-orbit
      computation (generation~2) with periodicity checking.
      \item \emph{MT2 + Periodicity + Perturb ST} (radio): MT2 reference orbit
      with a single-threaded perturbation pass for the high-resolution layer.
      \item \emph{MT2 + Periodicity + Perturb MT v1--v4} (radio entries):
      MT2 reference orbit with multithreaded perturbation.  Versions v1--v3
      represent different threading/scheduling strategies for the per-pixel
      perturbation pass; v4 is marked as broken.
      \item \emph{MT5 + Periodicity} (command): A newer multithreaded variant
      (generation~5) with periodicity checking.  Only enabled when perturbation
      support is available.

      \item \emph{GPU-Accelerated (see README)} (radio): Selects the
      GPU-accelerated perturbation path, which offloads high-precision
      reference-orbit computation to the GPU via the \code{HpSharkFloat}
      pipeline.  This option is a mode selection (radio button), not a one-shot
      action.  Intended for extremely deep zooms where CPU-based reference
      computation is prohibitively slow; may be slower than CPU at modest
      precision levels.  See \cref{subsec:ref-orbit-gpu} for details.

      \item \emph{Clear and Reload Reference Orbits} (command): Clears reference-orbit state and reloads it from persistent storage.
      \item \emph{Save Reference Orbits} (command): Saves reference-orbit state to persistent storage.
    \end{itemize}

  \item \emph{Palette Color Depth} (submenu): Controls palette selection,
    palette bit depth, and optional palette rotation.  Some entries may be
    enabled only when palette rotation support is available.
    \begin{itemize}
      \item \emph{Basic} (radio): Placeholder palette type; not currently
      generated.  Selecting it produces a black image.
      \item \emph{Default} (radio): Rainbow palette cycling through
      red, yellow, green, cyan, blue, magenta, and black.
      \item \emph{Patriotic} (radio): Red-and-blue palette.
      \item \emph{Summer} (radio): Red/green/cyan/white warm-tone palette.
      \item \emph{Random} (radio): Procedurally generated random palette;
      each selection produces the same random sequence (seed is fixed).
      Use \emph{Create Random Palette} to generate a fresh one.

      \item \emph{Create Random Palette} (command): Generates a new random palette.

      \item \emph{5-bit} (radio)
      \item \emph{6-bit} (radio)
      \item \emph{8-bit} (radio)
      \item \emph{12-bit} (radio)
      \item \emph{16-bit} (radio)
      \item \emph{20-bit} (radio)

      \item \emph{Palette Rotation} (command): Enters an interactive rotation
      mode.  The palette shifts by 10~units per frame while the image redraws
      continuously.  Moving the mouse more than 5~pixels from the starting
      position exits the mode and resets the palette to its original state.
      Useful for quickly previewing how different color offsets look.
    \end{itemize}

  \item \emph{Memory Management} (submenu): Controls memory-limit policy and
    whether reference orbit data is automatically saved.
    \begin{itemize}
      \item \emph{Enable Auto-Save Orbit (delete when done, default)} (radio):
      Automatically saves orbit data temporarily and deletes it when no longer
      needed.
      \item \emph{Enable Auto-Save Orbit (keep files)} (radio): Automatically
      saves orbit data and retains files on disk.
      \item \emph{Disable Auto-Save Orbit} (radio): Disables automatic orbit saving.

      \item \emph{Remove Memory Limits} (radio): Disables memory limiting behavior.
      \item \emph{Leave max of (1/2*ram, 8GB) free (default)} (radio): Constrains allocation to preserve system memory headroom.
    \end{itemize}

  \item \emph{Show Rendering Details} (command): Displays current rendering
  status and diagnostics, including position coordinates (real, imaginary,
  zoom level), the active render algorithm, and perturbation/LA state.
  The information is shown in a message box and also copied to the clipboard.

  \item \emph{Save} (submenu): Output and benchmarking commands.
    \begin{itemize}
      \item \emph{Save Location...} (command): Saves the current location/view parameters.
      \item \emph{Save High Res Bitmap} (command): Saves a high-resolution
      bitmap render.  The output resolution matches the current window
      dimensions (or a scaled multiple thereof).  A file-save dialog is
      presented to choose the output path.
      \item \emph{Save Iterations as Text} (command): Saves iteration counts in a text format.
      \item \emph{Save Bitmap Image} (command): Saves the current bitmap image.

      \item \emph{Save Reference Orbit as Text} (command): Saves the reference orbit in text form.
      \item \emph{Save Compressed Orbit as Text (simple)} (command)
      \item \emph{Save Compressed Orbit as Text (max)} (command)
      \item \emph{Save Compressed Orbit (Imagina/max)} (command): Saves the
      reference orbit in Imagina's file format with maximum compression,
      suitable for import into Imagina or later reload via
      \emph{Load~$\to$~Load Imagina Orbit}.

      \item \emph{Benchmark (5x, full recalc)} (command): Runs 5 full
      recalculation passes using all perturbation result types and reports
      overall and per-pixel rendering times in milliseconds.
      \item \emph{Benchmark (5x, intermediate only)} (command): Runs 5
      recalculation passes using only the medium-resolution perturbation
      results, which isolates intermediate-stage performance.

      \item \emph{Diff Imagina Orbits (choose two)} (command): Opens two
      file-chooser dialogs to select a pair of Imagina-format orbit files, then
      compares them element-by-element and reports any differences.  Useful for
      verifying that two orbit computations produce equivalent results.
    \end{itemize}

  \item \emph{Load} (submenu): Load view parameters and reference orbits.
    \begin{itemize}
      \item \emph{Load Location...} (command): Loads a previously saved location/view.
      \item \emph{Enter Location} (command): Manually enters a location/view
      via a dialog where the user enters the real part, imaginary part, and zoom level.
      \item \emph{Load Imagina Orbit (Match)...} (command): Loads an
      Imagina-format orbit file and matches it against the current view
      parameters (coordinates and precision) to validate compatibility.
      \item \emph{Load Imagina Orbit (Use Saved)...} (command): Loads a
      previously saved Imagina orbit file directly, using whatever parameters
      are stored in the file without matching against the current view.
    \end{itemize}

  Imagina is a separate Mandelbrot rendering application.  \FractalShark{}
  can import and export reference orbits in Imagina's file format, which is
  useful for cross-validating orbit computations between the two programs or
  for sharing precomputed orbits.

  \item \emph{Exit} (command): Exits the application.
\end{itemize}

\begin{figure}[!htbp]
  \centering
  \includegraphics[width=0.9\linewidth]{auto-renderer-artifact.png}
  \caption[Precision artifacts in HDRx32 GPU renderer]{Rendering artifacts arising from insufficient precision in the
  HDRx32 GPU renderer. Switching to HDRx64 or a CPU-based renderer resolves these
  artifacts.}
  \label{fig:auto-renderer-artifact}
\end{figure}

\subsection{Recalculation Commands}
\label{subsec:recalculation-commands}

Commands such as \emph{Recalculate} or \emph{Recalculate, Reuse Reference} trigger
an immediate re-render of the current view. These commands do not change any
persistent settings; they simply instruct the renderer to restart computation
using the current configuration. \emph{Recalculate, Reuse Reference} keeps the
existing high-precision reference orbit and only recomputes the per-pixel
perturbation pass, which is significantly faster when the view center has not
changed.

\subsection{Experimental and Debug Options}
\label{subsec:experimental-debug}

Some menu entries expose experimental features, debugging aids, or partially
implemented functionality. These options may be unstable, incomplete, or
subject to removal. \FractalShark{} is an experimental research project rather than
a polished end-user application. Users should expect occasional crashes,
rendering failures, or unresponsive behavior when exercising less common menu
options.  See also \cref{subsec:practical-usage}.

\subsection{Practical Usage Notes}
\label{subsec:practical-usage}

\begin{itemize}
  \item Long-running computations (renders, reference orbits, feature-finder
  searches) can be aborted by holding \code{Ctrl+Alt} for approximately three
  seconds.  An \code{AbortMonitor} background thread polls the key state every
  250\,ms and sets a stop flag once the debounce threshold is reached.  Pressing
  \code{Escape} cancels (resets) the abort flag, allowing computation to
  continue if it has not yet terminated.
  \item Some menu options may be temporarily disabled depending on the current
  render state or selected algorithm.
  \item If the application exits unexpectedly, a Windows minidump
  (\code{.dmp}) file may be generated in the same directory as the
  executable.  These dumps can be loaded in Visual Studio or WinDbg for
  post-mortem debugging.
\end{itemize}

Despite its rough edges, the popup menu provides a unified control surface for
exploring \FractalShark{}’s many rendering paths and numerical experiments, making
it the primary interface for both casual exploration and deep technical
investigation.


\section{Keyboard Shortcuts}
\label{sec:keyboard-shortcuts}

In addition to the popup menu, \FractalShark{} supports a set of single-key
shortcuts for navigation, recalculation, palette exploration, and tuning of
perturbation/linear-approximation parameters.  Most shortcuts operate relative
to the \emph{current mouse position} (converted into client coordinates at the
time the key is pressed), which makes them feel similar to context-menu actions.

Unless otherwise noted, shortcuts are \emph{case-sensitive} in the sense that
uppercase variants are invoked by holding \code{Shift}.  Many actions also
explicitly clear cached perturbation results prior to recomputation, trading
additional work for correctness or a ``clean'' run.

\paragraph{Navigation and view control.}
\begin{itemize}
  \item \code{z}: Zoom in by a fixed amount, centered at the mouse cursor.
  \item \code{Z} (\code{Shift+z}): Zoom out by a fixed amount, centered at the mouse cursor.
  \item \code{c}: Center the view at the mouse cursor (no forced reference recomputation).
  \item \code{C} (\code{Shift+c}): Center the view and clear all perturbation results first.
  \item \code{b}: Go back to the previous view.
  \item \code{r}: Square the current view (aspect normalization).
  \item \code{R} (\code{Shift+r}): Clear all perturbation results, then square the current view.
  \item Mouse \textbf{left-click and drag}: Zoom into a dragged rectangle.
        Holding \code{Shift} while dragging disables aspect-ratio enforcement.
\end{itemize}

\paragraph{Autozoom (experimental / buggy).}
\label{par:autozoom}

\begin{itemize}
  \item \code{a}: Autozoom using the ``Feature'' heuristic.  Forces the
        \code{GpuHDRx32PerturbedLAv2} render algorithm, centers on the
        perturbation reference point, computes a target zoom depth from the
        reference orbit's maximum radius, and zooms in incrementally while
        linearly interpolating the iteration count toward the reference orbit's
        target.
  \item \code{A} (\code{Shift+a}): Center the view at the mouse cursor, then
        autozoom using the ``Default'' heuristic.  Iteratively picks
        a zoom target via a weighted geometric mean of iteration counts
        (favoring high-iteration, center-proximate pixels) and zooms by a
        factor of~3 each step.
  \item \code{S} (\code{Shift+s}): Autozoom using the ``FilamentTip'' heuristic.
\end{itemize}
Holding \code{Ctrl+Alt} for approximately three seconds aborts either
autozoom variant (the abort is detected by the \code{AbortMonitor} thread;
pressing \code{Escape} cancels the abort and allows the zoom to resume).

\paragraph{Recalculation and benchmarking.}
Several keys force a recalculation (see also
\cref{subsec:recalculation-commands}); uppercase forms often clear cached
results first.
\begin{itemize}
  \item \code{i}: Force recalculation.
  \item \code{I} (\code{Shift+i}): Clear \emph{medium-res} perturbation results, then recalculate.
  \item \code{o}: Force recalculation.
  \item \code{O} (\code{Shift+o}): Clear \emph{all} perturbation results, then recalculate.
  \item \code{p}: Force recalculation.
  \item \code{P} (\code{Shift+p}): Clear \emph{LA-only} perturbation results, then recalculate.
\end{itemize}
(The in-application hotkey dialog groups \code{i/I}, \code{o/O},
\code{p/P}, and \code{r/R} under ``Recalculating and Benchmarking'' with
slightly different wording.  The authoritative behavior is the code path: these
keys primarily differ by which cached perturbation result class is cleared
before forcing recomputation.)

\paragraph{Reference compression tuning.}
These keys adjust compression error exponents and then repaint, typically after
clearing cached perturbation results.
\begin{itemize}
  \item \code{e}: Clear all perturbation results, reset compression error exponents to defaults,
        and repaint.
  \item \code{q}: Clear all perturbation results; \emph{increase} intermediate compression error
        exponent by 10 (more error, less memory).
  \item \code{Q} (\code{Shift+q}): Clear all perturbation results; \emph{decrease} intermediate compression
        error exponent by 10 (less error, more memory).
  \item \code{w}: Clear all perturbation results; \emph{increase} low-error compression exponent by 1.
  \item \code{W} (\code{Shift+w}): Clear all perturbation results; \emph{decrease} low-error compression exponent by 1.
\end{itemize}
\emph{Practical interpretation:} increasing the error exponent generally favors
memory reduction and speed at the expense of fidelity in reconstructed reference
segments; decreasing it pushes toward higher fidelity and higher memory use.

\paragraph{Linear approximation parameter tuning (powers-of-two adjustments).}
These keys adjust linear-approximation thresholds and period-detection thresholds.
After changing parameters, \FractalShark{} clears LA-only perturbation results and
forces a recalc.
\begin{itemize}
  \item \code{h}: Increase LA threshold scale exponents (less accurate / faster per-pixel).
  \item \code{H} (\code{Shift+h}): Decrease LA threshold scale exponents (more accurate / slower per-pixel).
  \item \code{j}: Increase LA period-detection threshold exponents.
  \item \code{J} (\code{Shift+j}): Decrease LA period-detection threshold exponents.
\end{itemize}

\paragraph{Palettes.}
\begin{itemize}
  \item \code{d}: Advance to the next palette lookup-table depth and redraw.
  \item \code{D} (\code{Shift+d}): Create a new random palette, select it, and redraw.
  \item \code{t}: Increase auxiliary palette depth (cycling auxiliary palette behavior) and redraw.
  \item \code{T} (\code{Shift+t}): Decrease auxiliary palette depth and redraw.
\end{itemize}

\paragraph{Iteration count.}
These keys scale the maximum iteration count aggressively for quick exploration:
\begin{itemize}
  \item \code{=} : Multiply max iterations by 24.
  \item \code{-} : Multiply max iterations by $2/3$.
\end{itemize}

\paragraph{Window manipulation and menu invocation.}
\begin{itemize}
  \item \textbf{Right-click}: Open the popup menu at the cursor position.
  \item \textbf{Shift+F10 / Menu key}: Also opens the popup menu (Windows context-menu invocation).
  \item \textbf{Alt + left-click drag} (windowed mode): Drag the window by treating the client area as a caption.
  \item \textbf{Mouse wheel}: Zoom in/out (scroll up zooms in; scroll down zooms out).
\end{itemize}

\paragraph{Abort and cancel.}
\begin{itemize}
  \item \textbf{Ctrl+Alt} (hold ${\sim}$3\,s): Abort the current long-running
        computation (render, reference orbit, feature-finder search).
        The \code{AbortMonitor} background thread polls every 250\,ms
        and sets the stop flag after 12~consecutive positive samples.
  \item \textbf{Escape}: Cancel (reset) the abort flag.  If pressed
        while the stop flag is set but computation has not yet
        terminated, allows the operation to resume.
\end{itemize}

\paragraph{Feature Finder.}
These keys invoke the periodic-point (feature) finder at the current mouse
position, or scan the entire viewport.  See \cref{sec:feature-finder}
for the algorithm details and \cref{sec:ui-popup-menu} for the
corresponding menu items under \emph{Navigate $\to$ Feature Finder}.
\begin{itemize}
  \item \code{n}: Find the nearest feature using direct iteration (single point at cursor).
  \item \code{N} (\code{Shift+n}): Scan the viewport on a $12\times12$ grid using direct iteration.
  \item \code{m}: Find the nearest feature using perturbation theory (single point at cursor).
  \item \code{M} (\code{Shift+m}): Scan the viewport on a $12\times12$ grid using perturbation theory.
  \item \code{,}: Find the nearest feature using linear approximation (single point at cursor).
  \item \code{<} (\code{Shift+,}): Scan the viewport on a $12\times12$ grid using linear approximation.
  \item \code{.}: Zoom to the most recently found feature (refines to high precision first if needed).
  \item \code{>} (\code{Shift+.}): Clear all found-feature overlays.
\end{itemize}

\paragraph{Responsiveness note.}
Many shortcuts trigger a recomputation or redraw immediately.  If \FractalShark{}
is executing a long-running render (especially a large GPU kernel), keyboard
input may appear delayed until the render completes.  In practice, the popup
menu and hotkeys are best used as \emph{configuration and relaunch} controls,
not as guaranteed real-time interaction during heavy computation.

\subsection{Commentary}

\begin{itemize}
\item An interesting experiment is running with LA~v2 + ``LA only''. This 
approach actually works pretty well once Linear Approximation kicks in at 
deeper zooms --- it provides a sense of what the actual image should look like 
but is very fast, since it does no perturbation. The images it produces are not 
precise, and often leave out the fine detail; however, the mode is enjoyable to play with 
when zooming in on a specific point.

\item The most interesting reference orbit calculation is at
\code{AddPerturbation\-ReferencePointMT3}. It includes a ``bad'' calculation
which is used for the ``scaled'' CUDA kernels. The multithreaded approach
handily beats the single-threaded implementation on my CPU in all scenarios.

\item There are CPU renderers, but they were mostly to learn/debug more easily, 
and are not optimized heavily. They are much easier to understand and reason 
about though.
\end{itemize}

% -----------------------------------------------------------------
\section{Automatic Zoom}
\label{subsec:autozoom}

The \code{AutoZoomer} class provides four heuristic strategies for
automated zoom animation.  Each heuristic repeats until a termination
condition is met (all pixels at the iteration limit, user abort, or
target reached).  The render pipeline depth of
$4 \times \text{NumRenderers}$ keeps several frames in flight for smooth
playback.

\begin{description}
  \item[Default (zoom divisor~3)]
    A weighted geometric mean of pixel iteration counts in a
    center-cropped region (borders trimmed to $\tfrac{1}{8}$ of each
    dimension).  Pixels whose iteration count exceeds the image average
    are weighted by their distance from the screen center and their
    iteration count relative to the observed maximum.  Each step zooms
    to $\tfrac{1}{3}$ of the current view, centered on the weighted
    target.

  \item[Max (zoom divisor~32)]
    The first pixel that attains the maximum iteration count is selected,
    and the view zooms aggressively---32$\times$ per step.  The sequence
    terminates when every pixel reaches the iteration limit.

  \item[FilamentTip (zoom divisor~8)]
    A multi-stage scoring function targets isolated filament tips.  For
    each candidate pixel whose iteration count exceeds the image
    average, the heuristic samples eight compass directions at a radius
    of 12~pixels and counts high-iteration neighbors.  A composite
    tip score combines isolation (fewer high neighbors is preferable),
    elevation above the average, and distance from the screen center.
    The pixel with the highest score becomes the zoom target.

  \item[Feature (zoom divisor~3)]
    A GPU-driven strategy that forces \code{GpuHDRx32PerturbedLAv2},
    locates the nearest periodic point via the feature finder
    (\cref{sec:feature-finder}), and pre-computes the reference orbit
    at the target zoom depth.  An animated sequence is then generated
    with a $1.1\times$ zoom factor per frame.  Iteration counts are
    linearly interpolated from the current value to the target across
    all frames.
\end{description}

% -----------------------------------------------------------------
\section{FractalTray Batch Renderer}
\label{subsec:fractaltray}

\code{FractalTray} is a companion utility that renders zoom-sequence
animations in batch, writing each frame to a network share or local
directory as a BMP file.  The application runs as a system-tray icon
and is designed for unattended overnight rendering.

\subsubsection*{Coordinate format}

Source and destination coordinates share a single line format:

\begin{lstlisting}
width height minX minY maxX maxY iters gpuAntialiasing
\end{lstlisting}

\noindent
The coordinate values (\code{minX}, \code{minY}, \code{maxX},
\code{maxY}) are arbitrary-precision decimal strings; the remaining
fields are integers.  Width and height specify the output image
dimensions in pixels.

\subsubsection*{Generating a location file}

The dialog provides two text fields---source coordinates and
destination coordinates---together with a scale factor and output
resolution.  Clicking \emph{Generate} computes the complete frame
sequence and writes it to \code{locations.txt}.

Frame generation proceeds as follows.  Let $S = (S_\mathrm{minX},
S_\mathrm{minY}, S_\mathrm{maxX}, S_\mathrm{maxY})$ and $D$
denote the source and destination bounding boxes respectively, and
let $f$ denote the scale factor (default~75).  Starting from $C = S$,
each frame updates the current bounding box:
\[
  C \;\leftarrow\; C + \frac{D - C}{f}
\]
This geometric approach progressively halves the remaining distance
to the destination at a rate governed by $f$.  Larger values of $f$
produce a slower, smoother zoom with more frames; smaller values
produce a faster zoom with fewer frames.  The iteration count is
linearly interpolated from the source to the destination value
across the total number of frames.

The frame count is determined by repeating the zoom step until the
destination bounding-box area exceeds $90\%$ of the current area,
at which point the current and destination views are effectively
identical.

Frames are written to \code{locations.txt} in \emph{reverse}
order---deepest zoom first---so that during rendering, the first
frame computed is the deepest.  This ordering allows the reference
orbit computed for the deepest frame to be reused for subsequent
(wider) frames, avoiding redundant orbit computation.

\subsubsection*{Rendering}

On startup, \code{FractalTray} checks for an existing
\code{locations.txt}.  If found, rendering begins immediately and
the dialog minimizes to the system tray, where the tray icon changes
to indicate active rendering.

For each line in the location file, the renderer:

\begin{enumerate}
  \item Parses the coordinate string and output filename.
  \item Checks whether the output file already exists (BMP or PNG) in
        the target directory; if so, the frame is skipped.  This
        skip-if-exists behavior allows the process to be interrupted
        and restarted without re-rendering completed frames.
  \item Computes the required MPIR precision from the bounding box
        using \code{PrecisionCalculator}.
  \item Configures the \code{Fractal} object with the frame's
        dimensions, iteration count, precision, coordinates, and
        palette.
  \item Calls \code{CalcFractal(true)} (the direct synchronous path),
        which writes iteration data to \code{m\_CurIters}.
  \item Saves the result via \code{SaveCurrentFractal()}.
\end{enumerate}

The rendering loop runs on a background \code{std::jthread}.
Cancellation is cooperative: the tray icon's \emph{Exit} command
requests a stop via the thread's stop token, and the loop checks
the token at the start of each frame.

\subsubsection*{Configuration}

\begin{description}
  \item[Scale factor] Controls zoom speed.  The default value of
    $75$ produces a gradual, cinematic zoom; lower values (e.g.~10)
    zoom faster with fewer frames.
  \item[Resolution] Output width and height in pixels.  The default
    is $3840 \times 1600$.
  \item[Output directory] Controlled by the \code{FilePrefix}
    constant, currently a UNC network path.  Each frame is saved as
    \code{<prefix>output-NNNNNN.bmp}.
\end{description}

% -----------------------------------------------------------------
\section{Functional Test Suite}
\label{subsec:crummy-test}

\code{CrummyTest} is a self-contained functional test suite invoked
from the application's right-click context menu
(\code{IDM\_BASICTEST}, labeled ``Run Basic Test'').  The suite
exercises the direct rendering path
(\cref{subsec:two-render-paths}): each test calls \code{Drain()}
first to ensure no pool workers remain active, then invokes
\code{CalcFractal(true)}, which writes results to \code{m\_CurIters}.

The top-level \code{TestAll()} method runs the full suite in sequence:
\begin{itemize}
  \item \code{TestWindowResize}---resizes the application window and
        verifies that a re-render completes correctly.
  \item \code{TestBasic}---iterates through every
        \code{RenderAlgorithm} entry, rendering the default view,
        View~5, and View~11 with both 32-bit and 64-bit iteration
        types.
  \item \code{TestReferenceSave}---saves and reloads reference orbits,
        verifying round-trip correctness.
  \item \code{TestVariedCompression}---tests reference-orbit
        compression at varying error thresholds.
  \item \code{TestImaginaLoad}---loads orbits stored in the Imagina
        format (\cref{subsec:imagina-format}).
  \item \code{TestStringConversion}---validates string-to-precision
        and precision-to-string conversion utilities.
  \item \code{TestPerturbedPerturb}---exercises perturbed iteration
        chains for numerical correctness.
  \item \code{TestGrowableVector}---validates the
        \code{GrowableVector} data structure
        (\cref{sec:disk_backed_growable_vectors}).
\end{itemize}
A separate \code{TestReallyHardView27} (not included in
\code{TestAll}) renders View~\#27, a period ${\sim}28$~billion point
requiring hours of computation; the test serves as an overnight
stress benchmark rather than a routine regression check.

% -----------------------------------------------------------------
\section{Kernel List}
\label{sec:kernel-list}

This section simply lists all the CUDA kernels in \FractalShark{}.

\begin{itemize}

\item \textbf{Mandelbrot base kernels} (\cref{sec:mandel-base-kernels})
  \begin{itemize}
    \item \code{mandel\_1x\_float}
    \item \code{mandel\_1x\_double}
    \item \code{mandel\_2x\_float}
    \item \code{mandel\_2x\_double}
    \item \code{mandel\_4x\_float}
    \item \code{mandel\_4x\_double}
    \item \code{mandel\_hdr\_float}
  \end{itemize}

\item \textbf{Mandelbrot perturbation kernels}
  \begin{itemize}
    \item \code{mandel\_1x\_double\_perturb\_bla}
    \item \code{mandel\_1xHDR\_float\_perturb\_bla} (\cref{sec:bilinear-approx})
    \item \code{mandel\_1xHDR\_float\_perturb\_lav2} (\cref{sec:la-v2-perturb,sec:perturb-only})
  \end{itemize}

\item \textbf{High-precision reference kernels} (\cref{subsec:ref-orbit-gpu})
  \begin{itemize}
    \item \code{HpSharkReferenceGpuKernel}
    \item \code{HpSharkReferenceGpuLoop}
  \end{itemize}

\item \textbf{Utility kernels}
  \begin{itemize}
    \item \code{antialiasing\_kernel}
    \item \code{max\_reduce}
    \item \code{max\_kernel}
  \end{itemize}

\item \textbf{Addition kernels} (\cref{sec:hp-add})
  \begin{itemize}
    \item \code{AddKernel}
    \item \code{AddKernelTestLoop}
  \end{itemize}

\item \textbf{Multiplication (NTT) kernels} (\cref{sec:ntt-multiply})
  \begin{itemize}
    \item \code{MultiplyKernelNTT}
    \item \code{MultiplyKernelNTTTestLoop}
  \end{itemize}

\item \textbf{Disabled kernels}
  \begin{itemize}
    \item \code{mandel\_2x\_float\_perturb\_setup}
    \item \code{mandel\_2x\_float\_perturb}
  \end{itemize}

\end{itemize}

\subsection{Commentary}

\begin{itemize}

\item All the per-pixel CUDA kernels are in \code{gpu\_render.cu}. The two
most interesting are probably \code{mandel\_\allowbreak 1xHDR\_\allowbreak float\_\allowbreak perturb\_bla} and
\code{mandel\_\allowbreak 1xHDR\_\allowbreak float\_\allowbreak perturb\_lav2}, which are the ones I've spent the
most time on lately. For better performance at low zoom levels, consider
\code{mandel\_1x\_float\_perturb}, which leaves out linear approximation
and just does straight perturbation up to $\sim 10^{30}$, which corresponds
with the 32-bit float exponent range.

\item The \code{mandel\_1x\_float} (\cref{sec:mandel-1x-float}) is the classic 32-bit float Mandelbrot and 
is screaming fast on a GPU. This one is optimized with fused multiply-add for 
fun even though it is of limited practical use because one can barely zoom in before
pixellation sets in.

\end{itemize}

